\section{问题二的模型与求解}

本文采用\textbf{期望定位损失}与{完成概率}评价第二次检测的效果，分别考察检测后还剩下多少位置不确定性，以及无需补测即可确定保证清除位置的可能性。

其中，“定位损失”用剩余可行区域的最小包围圆半径衡量尚未消除的位置不确定性；若整个区域已位于当前检测点的20米以内，能够原地完成光学精确定位，则损失取零。将各种反馈下的损失按概率求平均，就得到期望定位损失，数值越小越好。“完成概率”则表示第二次检测后，剩余区域的包围圆半径在20米内的概率；此时无需继续无线电测向，移动到该圆圆心即可完成光学定位与清除。

需要注意的是，源位置、接收半径和各地点的测角误差均为固定未知量。本文采用贝叶斯方法，以\textbf{主观概率}描述在已有信息下对这些未知量取值的认识。获得新观测后，通过贝叶斯公式更新其概率分布，改变的是我们对未知量的判断，其实际取值保持不变，因此与题目规定的同一地点误差固定相容。

在此基础上，本文建立\textbf{基于贝叶斯更新的加权选点模型}。先根据首次观测更新对固定未知量的判断，再预测第二次反馈，计算上述两个指标，并构造定位效果综合评价函数。权重确定后，最小化该函数即可作出选点决策；改变权重，则得到不同定位偏好下的一组候选位置。

\subsection{准备工作：由观测信息构造定位区域}

设第二个检测点为 \(q\)，它的位置是本问需要决定的量。为简化坐标表示，利用圆的旋转对称性，不妨将首次检测点写成 \(s_1=(a,0)\)，首次示向度记为 \(\beta\)，其中 \(a\ge0\)、\(\beta\in[-\pi,\pi)\) 均为确定的参数。若在 \(q\) 处得到正常示向度 \(\theta\)，则令 \(s_2=q\)、\(\theta_1=\beta\)、\(\theta_2=\theta\)，并沿用问题一的角域记号 \(W_1,W_2\)。源位置 \(g\) 到两次检测点的距离分别简记为 \(d_1(g)=\|g-s_1\|\)、\(d(g;q)=\|g-q\|\)。

首次测得正常示向度，说明源位于对应角域内，能够被接收，并且尚未近到触发强信号。因此，首次观测后的源位置可行区域为 \(K_1=\{g\in\Omega\cap W_1:r_{\mathrm{strong}}<d_1(g)\le1500\}\)。对于其中一个可能位置 \(g\)，接收半径还必须满足 \(\max\{1000,d_1(g)\}\le R\le1500\)。

在 \(q\) 处进行第二次检测，可能得到强信号（源距\(q\)不超过5米）、无信号（源距大于接收半径）或正常示向度，分别记为 \(\mathrm H\)、\(\mathrm N\) 和 \(\theta\)。
在这里，无信号也能提供信息，这是由于首次探测已成功接收，且两次检测对应同一个固定接收半径，第二次无信号表明源距 \(q\) 比源距 \(s_1\) 更远，同时超过1000米。将这些新信息与 \(K_1\) 取交集，得到三种反馈后的剩余区域

\begin{equation}
\begin{aligned}
K_{\mathrm H}(q)&=\{g\in K_1:d(g;q)\le r_{\mathrm{strong}}\},\\
K_{\mathrm N}(q)&=\{g\in K_1:d(g;q)>\max\{1000,d_1(g)\}\},\\
K_\theta(q)&=\{g\in K_1\cap W_2:r_{\mathrm{strong}}<d(g;q)\le1500\}.
\end{aligned}
\label{eq:q2-pair-regions}
\end{equation}

其中，\(W_2\) 由检测点 \(q\) 和读数 \(\theta\) 确定。若某种反馈对应的区域为空，它就不可能在已有观测条件下出现，后续不计入评价。

沿用问题一的方法，用最小包围圆半径 \(\rho(K)\) 衡量区域 \(K\) 还有多大，半径越小，源的位置就确定得越精细。第二次反馈为 \(z\) 时，将剩余区域 \(K_z(q)\) 的这一半径简记为 \(\rho_z(q)\)。

若第二个检测点距 \(K_1\) 超过1500米，则必然只能得到无信号，因而无法缩小已有区域。因此，将距 \(K_1\) 不超过1500米、且不同于首次检测点的位置作为搜索范围，记为 \(Q\)（允许位于目标圆域 \(\Omega\) 外）。

\subsection{三种反馈下的定位损失}

按照前述定义，定位损失以米为单位，数值越小，表示检测后的定位效果越好。根据第二次检测的反馈，分别按以下规则计算。

\begin{enumerate}[label=(\arabic*),leftmargin=2.2em]
  \item \textbf{强信号。} 源距当前检测点不超过5米，已在20米的光学定位范围内，因此可以原地确定源的准确位置，取 \(L_{\mathrm H}(q)=0\)。
  \item \textbf{无信号。} 源距当前检测点超过1000米，无法原地完成光学定位。因此，以无信号反馈后剩余区域的最小包围圆半径作为损失，取 \(L_{\mathrm N}(q)=\rho_{\mathrm N}(q)\)。
  \item \textbf{正常示向度。} 根据新读数 \(\theta\) 更新定位区域，再判断该区域是否整体位于当前检测点的20米范围内。若是，则可以原地完成光学定位，取 \(L_\theta(q)=0\)；否则，以更新后区域的最小包围圆半径作为损失，取 \(L_\theta(q)=\rho_\theta(q)\)。
\end{enumerate}

\subsection{基于贝叶斯更新的反馈预测}

对于同一个检测点，不同反馈留下的区域大小不同，能否达到清除条件也不同。为了计算期望定位损失和完成概率，需要先用首次观测更新对源位置和接收半径的判断，再预测第二次各种反馈出现的可能性。

\subsubsection{未知量的先验描述}

按照前文模型假设，源位置 \(G\) 在 \(\Omega\) 内按面积均匀分布，接收半径 \(R\sim\operatorname{U}[1000,1500]\)，两个不同检测点的误差 \(E_1,E_2\sim\operatorname{U}[-\alpha,\alpha]\)，观测前认为这些未知量相互独立。均匀性和独立性是为了计算反馈概率而增加的假设，不能仅由题目给出的误差界与接收范围推出。后续的期望定位损失与完成概率均基于上述先验假设和观测信息计算，其结果反映这些条件下的预测评价。

同一地点的多次检测始终共享同一个误差，不能重新抽样，也不能靠重复测量取平均消除误差。本模型只预测不同于首次检测点的新位置。对于观测后产生的关联，例如较远的源位置必须对应较大的接收半径，后续计算会一并保留。

\subsubsection{首次观测后的概率更新}

将首次正常示向度信息记为 \(I_1=(s_1,\beta)\)。可行区域 \(K_1\) 告诉我们哪些位置没有被排除，但其中的位置未必同样可能。例如，源距不超过1000米时，题目允许的任何接收半径都能解释首次接收；距离超过1000米后，只有较大的接收半径才能解释这一观测。因此，应根据接收条件调整不同位置的概率。

贝叶斯更新将原有概率与观测的似然相乘，再归一化。似然表示假定某个源位置成立时，得到已有观测的可能程度。在 \(K_1\) 内，均匀测角误差对应的角度似然相同，位置权重的差别来自能否接收到信号。距离为 \(d\) 时，这一接收概率为

\begin{equation}
 p_{\mathrm{rec}}(d)=\Pr(R\ge d)=
 \begin{cases}
 1,&d\le1000,\\
 (1500-d)/500,&1000<d<1500,\\
 0,&d\ge1500.
 \end{cases}
\label{eq:q2-pair-weight}
\end{equation}

将上述权重用于 \(K_1\) 内的各个位置，并使总概率等于1，得到首次观测后的源位置密度

\begin{equation}
 p_1(g)=\frac{\mathbf1_{K_1}(g)p_{\mathrm{rec}}(d_1(g))}{C_1},
 \qquad C_1=\int_{K_1}p_{\mathrm{rec}}(d_1(g))\,\mathrm dg>0.
\label{eq:q2-pair-posterior}
\end{equation}

其中 \(C_1\) 是归一化系数。该结果说明，距离超过1000米后，越远的位置得到的权重越低，因为能解释首次接收的半径取值越来越少。

图~\ref{fig:q2-posterior} 固定首次检测点为 \(s_1=(1200,0)\) 米，比较 \(\beta=30^\circ\) 与 \(90^\circ\) 时的后验分布。示向度决定保留哪个角域，而圆域边界进一步限制源的位置：\(30^\circ\) 方向的角域较早遇到边界，后验集中在较小的区域；\(90^\circ\) 方向保留的距离范围更长，其中超过1000米的部分还因接收概率下降而降低权重。这说明，当圆域边界参与截断时，示向度会改变后验分布的形状与密度大小。两幅热力图采用同一色标，并按源距与相对示向角展开显示；颜色仍表示单位面积上的位置概率密度。

接收半径本身也需要更新。若源位于 \(g\)，首次接收已经排除了小于 \(d_1(g)\) 的半径。因此，对于后验密度为正的位置，有 \(R\mid G=g,I_1\sim\operatorname{U}[\max\{1000,d_1(g)\},1500]\)。预测第二次检测时，要使用这个更新后的范围，不能再次把 \(R\) 当作与源位置无关的均匀量。

\begin{figure}[htbp]
  \centering
  \includegraphics[width=160mm]{q2_posterior_general.pdf}
  \caption{非圆心检测点处，不同首次示向度对应的后验位置密度}
  \label{fig:q2-posterior}
\end{figure}

\subsubsection{第二次反馈的概率预测}

现在固定一个待评价的检测点 \(q\)，考虑源的每个可能位置 \(g\)。因为接收半径固定，两次都能接收要求半径至少达到两次源距中较大的一个。记这个距离为 \(d_\vee(g;q)=\max\{d_1(g),d(g;q)\}\)，则两次均能接收的概率权重为 \(p_{\mathrm{rec}}(d_\vee)\)。从首次能够接收的权重 \(p_{\mathrm{rec}}(d_1)\) 中减去它，就得到首次接收、第二次无信号的权重。

再结合三种反馈各自的位置约束，得到用于计算反馈概率的权重函数

\begin{equation}
\begin{aligned}
 b_{\mathrm H}(g;q)&=\mathbf1_{K_{\mathrm H}(q)}(g)p_{\mathrm{rec}}(d_1(g)),\\
 b_{\mathrm N}(g;q)&=\mathbf1_{K_1}(g)\bigl[p_{\mathrm{rec}}(d_1(g))-p_{\mathrm{rec}}(d_\vee(g;q))\bigr],\\
 b_\theta(g;q)&=\frac{\mathbf1_{K_\theta(q)}(g)}{2\alpha}p_{\mathrm{rec}}(d_\vee(g;q)).
\end{aligned}
\label{eq:q2-pair-feedback-weights}
\end{equation}

其中，正常示向度的读数还受新地点测角误差影响。按照不同地点误差独立且均匀的假设，允许的角度范围内密度为 \(1/(2\alpha)\)，因此第三项包含这个因子。将各个可能源位置的权重积分，并除以首次观测的归一化系数，得到

\begin{equation}
\begin{aligned}
 P_{\mathrm H}(q)&=\frac{1}{C_1}\int_{K_1}b_{\mathrm H}(g;q)\,\mathrm dg,
 &P_{\mathrm N}(q)&=\frac{1}{C_1}\int_{K_1}b_{\mathrm N}(g;q)\,\mathrm dg,\\
 f_\Theta(\theta;q)&=\frac{1}{C_1}\int_{K_1}b_\theta(g;q)\,\mathrm dg.
\end{aligned}
\label{eq:q2-pair-prediction}
\end{equation}

\(P_{\mathrm H}(q)\) 和 \(P_{\mathrm N}(q)\) 分别是强信号与无信号出现的概率。正常示向度可以连续变化，因此用密度 \(f_\Theta(\theta;q)\) 描述；将它在某个角度区间内积分，就得到读数落入该区间的概率。在全部角度上积分所得的正常反馈概率，与前两种概率相加等于1。

\subsection{定位效果的综合评价与选点决策}

至此，对于每个待选位置 \(q\)，我们既知道不同反馈会留下多大的定位损失，也知道这些反馈出现的概率。以 \(Z\) 表示尚未得到的第二次反馈，将损失按预测概率求平均，得到期望定位损失

\begin{equation}
 J(q)=\mathbb E[L_Z(q)\mid I_1,q]
 =P_{\mathrm N}(q)\rho_{\mathrm N}(q)
 +\int_0^{2\pi}f_\Theta(\theta;q)L_\theta(q)\,\mathrm d\theta.
\label{eq:q2-pair-expected-objective}
\end{equation}

右端第一项对应无信号的结果，第二项将各种正常示向度对应的损失按概率密度积分。强信号能够直接完成原地定位，损失为零，所以不再单列一项。

另一方面，由问题一的结论，若第二次反馈后满足 \(\rho_Z(q)\le20\) 米，就已经能够确定一个保证清除的位置，不再需要补测。这里允许移动到剩余区域的最小包围圆圆心，与前文“在当前检测点原地完成光学定位”的条件不同。将这一事件的概率记为 \(P_{\mathrm{finish}}(q)\)，则

\begin{equation}
\begin{aligned}
 P_{\mathrm{finish}}(q)
 &=\Pr\bigl(\rho_Z(q)\le20\mid I_1,q\bigr)\\
 &=P_{\mathrm H}(q)+P_{\mathrm N}(q)\mathbf1_{\{\rho_{\mathrm N}(q)\le20\}}\\
 &\quad+\int_0^{2\pi}f_\Theta(\theta;q)\mathbf1_{\{\rho_\theta(q)\le20\}}\,\mathrm d\theta.
\end{aligned}
\label{eq:q2-completion-probability}
\end{equation}
其中，\(\mathbf1_{\{\cdot\}}\) 在条件满足时取1，否则取0。强信号一定满足完成条件；无信号和正常示向度则都要检查剩余区域能否被20米圆覆盖。无信号虽然不能原地定位，但也可能将可行区域缩小到足以确定清除位置。

图~\ref{fig:q2-feedback-regions} 具体说明了20米包围圆条件。这里取首次点为原点、首次示向正东，第二检测点为 \(q=(813.82,521.02)\) 米。若新读数为 \(270^\circ\)，剩余区域的包围圆半径约为16.87米，移动到其圆心即可保证清除，无需补测；若读数为 \(300^\circ\)，半径约为30.43米，任何20米圆都无法覆盖整个区域，因而仍需补测。两幅局部图使用相同的长度尺度，虚线圆的半径均为20米。

\begin{figure}[htbp]
  \centering
  \includegraphics[width=130mm]{q2_feedback_regions.pdf}
  \caption{两种示向度反馈下的剩余区域与完成条件}
  \label{fig:q2-feedback-regions}
\end{figure}

上述两个指标分别反映平均剩余不确定性和本次检测后结束测向的可能性。例如，在首次点位于原点、示向正东的算例中，将完成概率纳入权衡后，数值优选点的期望损失可由25.08米增至25.67米，而完成概率由26.62\%提高至37.34\%。这说明\textbf{较小的平均损失增量可能换来明显更高的完成概率}，因此有必要同时考虑两个指标。

最后，以任务给定的20米光学定位距离作为尺度，并将 \(J(q)\) 除以20，消除其与概率之间的量纲差异。设 \(w\in[0,1]\) 为完成概率指标的权重，构造\textbf{定位效果综合评价函数} \(F_w(q)\)，并将第二检测点的选择写为

\begin{equation}
 \min_{q\in Q}\quad F_w(q)
 =(1-w)\frac{J(q)}{20}+w\bigl[1-P_{\mathrm{finish}}(q)\bigr].
\label{eq:q2-weighted-objective}
\end{equation}

\(F_w\) 越小，综合定位效果越好。\(w=0\) 时，只最小化期望定位损失；\(w=1\) 时，只最大化完成概率；中间取值表示两者之间的权衡。权重反映对定位效果的偏好，不是题目给定的物理参数。固定 \(w\) 后，求解式~\eqref{eq:q2-weighted-objective} 即得到该偏好下的选点决策。

沿用图~\ref{fig:q2-feedback-regions} 的原点算例，该检测点是 \(w=0.5\) 时得到的数值优选点。图~\ref{fig:q2-feedback} 给出其第二次示向度的预测密度，并用阴影标出满足完成条件的反馈。本例强信号与无信号的概率均为0，因此阴影面积就是完成概率，约为37.34\%。密度较高的读数未必能使区域缩小到20米圆内，这说明选点时需要同时考虑反馈出现的可能性及其留下的区域大小。

\begin{figure}[htbp]
  \centering
  \includegraphics[width=150mm]{q2_feedback_distribution.pdf}
  \caption{第二次示向度的预测密度及满足完成条件的反馈}
  \label{fig:q2-feedback}
\end{figure}

\subsection{由权重变化确定候选区域}

题目没有指定上述两项指标的相对重要程度，因此不必预先限定唯一的权重。让 \(w\) 在0至1之间变化，汇集各个权重下的最优位置，便得到第二检测点的候选区域

\begin{equation}
 Q_{\mathrm{cand}}=\bigcup_{w\in[0,1]}\arg\min_{q\in Q}F_w(q).
\label{eq:q2-candidate-region}
\end{equation}

每个入选位置都对应至少一种定位偏好。使用时先确定权重，再选取该权重对应的位置即可。若出现并列最优点，则保留这些等效选择；例如，首次点位于圆心时，关于首次示向方向对称的位置具有相同评价。在只重视完成概率且出现并列解时，可优先报告期望损失较小的代表点。

实际计算时，对若干权重分别求解并汇集优选点。候选区域可能表现为曲线上的点集或若干分离部分，不一定具有非零面积；具体求解方法与结果如下。

\subsection{数值求解方法}

对一般的 \((a,\beta)\)，反馈区域随检测点改变，评价函数同时包含分段概率积分和最小包围圆计算，难以写成统一、简洁的显式选点公式。因此，本文采用\textbf{确定性数值积分与网格、多初值局部搜索相结合的方法}，输入首次观测及权重，输出检测点和两项评价指标。

\textbf{第一步：构造区域与搜索范围。} 由 \((a,\beta)\) 计算 \(K_1\) 和归一化系数 \(C_1\)，检查首次观测是否可行。计算时以首次点为临时原点、首次示向为横轴，便于处理狭窄角域；最终将选点坐标平移、旋转回题目的坐标系。若首次区域已满足20米包围圆条件，则直接报告无需再次测向。

\textbf{第二步：计算一个检测点的评价值。} 固定 \(q\)，分别构造强信号、无信号和各正常读数对应的区域，用边界交点与圆弧计算最小包围圆。位置积分在极坐标下进行，包含面积因子 \(r\)；在接收权重和区域约束变化处划分积分区间，再用高斯求积计算反馈概率。对正常示向度的外层积分，还要在区域形状变化及包围圆半径穿过20米的位置分段，并自适应细分，从而分别累加 \(J(q)\) 与 \(P_{\mathrm{finish}}(q)\)。这种确定性计算避免了随机抽样波动干扰相近检测点的比较。

\textbf{第三步：搜索优选位置。} 先在搜索范围的外包矩形内布置100米网格，再在可行区域附近补充较密的网格与沿示向方向的采样点。每次评价均检查 \(q\ne s_1\) 且距 \(K_1\) 不超过1500米，不合格的位置不参与比较。对固定权重，选取5个位置相互分开的较优网格点，并加入初始区域圆心或已有优选点作为参考起点，分别进行 Nelder--Mead 单纯形搜索。该方法只比较函数值，适合本模型中导数不易稳定计算的情形；随后对两个较优结果提高积分精度并再次细化。局部细化以坐标收缩至厘米量级、目标值极差不超过 \(2\times10^{-9}\) 为停止条件，最多迭代130次。同一点的两项指标可供不同权重复用。

\textbf{第四步：复算与核验。} 最终候选点采用更高阶求积和更密的阈值扫描复算，并检查0.1、1、5米邻域内是否存在更优点。新增两组情形的期望损失在提高精度前后变化小于 \(2\times10^{-6}\) 米，概率归一化残差小于 \(10^{-7}\)。另以独立射线交区间和支撑圆计算核查区域几何，并用每组20000次联合先验抽样核对期望损失与完成概率。数值搜索与这些检查提供的是优选解及其计算依据，不构成连续空间中的全局最优证明。

\subsection{求解结果与第二检测点的选择}

\subsubsection{不同首次观测下的典型结果}

为分辨首次位置与示向度的影响，表~\ref{tab:q2-typical-results} 统一取 \(w=0.5\)，比较圆心、非圆心的不同朝向及靠近边界的情形。表中均使用以目标圆心为原点的坐标，存在等效选点时列出一个代表。

\begin{table}[htbp]
  \centering
  \caption{典型首次观测下的数值优选点（\(w=0.5\)）}
  \label{tab:q2-typical-results}
  \small
  \renewcommand{\arraystretch}{1.08}
  \begin{tabular}{rrrrr}
    \toprule
    \(a\) / 米 & \(\beta\) / 度 & 第二检测点 \(q\) / 米 & \(J(q)\) / 米 & \(P_{\mathrm{finish}}(q)\) / \% \\
    \midrule
    0 & 0 & \((813.8,\,521.0)\) & 25.67 & 37.34 \\
    1200 & 30 & \((1459.2,\,467.7)\) & 12.29 & 100.00 \\
    1200 & 90 & \((695.7,\,798.1)\) & 24.12 & 40.38 \\
    1200 & 150 & \((755.7,\,858.1)\) & 25.67 & 37.34 \\
    1700 & 0 & \((1770.1,\,42.0)\) & 1.88 & 100.00 \\
    \bottomrule
  \end{tabular}
\end{table}

表~\ref{tab:q2-typical-results} 说明，首次点到圆心的距离相同，并不必然对应相近的定位效果。例如，\(a=1200\) 米时，\(30^\circ\) 方向的初始区域受圆域边界截断较多，第二次检测的完成概率可达100\%；\(90^\circ\) 方向保留范围更长，完成概率约为40.38\%。\(150^\circ\) 方向的首次角域在接收范围内完整落入圆域，其定位几何与圆心算例等价，因此平移、旋转选点后得到相同的两项指标。靠近边界的 \(a=1700\)、\(\beta=0\) 情形则具有更小的剩余不确定性。

表中完成概率为100\%而期望损失仍为正，是因为检测点未必能够原地光学定位；允许随后移动到剩余区域的圆心时，则可无需补测完成定位。上述概率均是给定先验与首次观测下的预测值。

\subsubsection{首次位于原点、示向正东时的权重比较}

当 \(a=0,\beta=0\) 时，关于横轴对称的两个检测点具有相同评价。表~\ref{tab:q2-origin-weights} 列出11个权重对应的数值优选点，以 \(\pm\) 同时表示两支对称候选位置。

\begin{table}[htbp]
  \centering
  \caption{原点、正东算例中不同权重的选点与定位效果}
  \label{tab:q2-origin-weights}
  \small
  \renewcommand{\arraystretch}{1.08}
  \begin{tabular}{rrrr}
    \toprule
    权重 \(w\) & 第二检测点 \(q\) / 米 & \(J(q)\) / 米 & \(P_{\mathrm{finish}}(q)\) / \% \\
    \midrule
    0 & \((906.6,\,\pm572.8)\) & 25.08 & 26.62 \\
    0.05 & \((895.8,\,\pm562.3)\) & 25.09 & 28.77 \\
    0.1 & \((885.9,\,\pm554.0)\) & 25.11 & 30.40 \\
    0.2 & \((867.4,\,\pm541.8)\) & 25.20 & 32.84 \\
    0.35 & \((840.9,\,\pm529.6)\) & 25.39 & 35.42 \\
    0.5 & \((813.8,\,\pm521.0)\) & 25.67 & 37.34 \\
    0.65 & \((785.1,\,\pm513.1)\) & 26.08 & 38.81 \\
    0.8 & \((754.0,\,\pm502.4)\) & 26.65 & 39.89 \\
    0.9 & \((731.4,\,\pm490.7)\) & 27.20 & 40.37 \\
    0.95 & \((719.1,\,\pm482.2)\) & 27.56 & 40.51 \\
    1 & \((705.5,\,\pm470.7)\) & 28.02 & 40.57 \\
    \bottomrule
  \end{tabular}
\end{table}

由表~\ref{tab:q2-origin-weights}，只考虑期望损失时，优选点约为 \((906.6,\pm572.8)\) 米。两项归一化指标等权时，选取 \(w=0.5\)，得到 \(q\approx(813.8,\pm521.0)\) 米；相较 \(w=0\)，期望损失增加约0.60米，完成概率提高约10.72个百分点。若只追求完成概率，则取 \(w=1\)，优选点约为 \((705.5,\pm470.7)\) 米，完成概率约为40.57\%，相应期望损失增至28.02米。

图~\ref{fig:q2-candidates} 将这些结果画在位置平面和指标平面上：左图显示两支镜像候选点，二者之间的位置不属于已计算的候选点；右图的曲线逐渐变平，说明继续提高完成概率需要付出更多期望损失。由此，本问的输出是随权重变化的一组候选位置；确定定位偏好后，采用对应权重的选点。对于表外的首次观测，按前述数值流程重新求解即可。

\begin{figure}[htbp]
  \centering
  \includegraphics[width=160mm]{q2_candidates.pdf}
  \caption{11个权重下的候选检测点及评价指标（\(a=0,\beta=0\)）。颜色表示完成概率权重，虚线连接相邻权重的数值结果。}
  \label{fig:q2-candidates}
\end{figure}

\FloatBarrier
